\documentclass[12pt,a4paper]{article}
\usepackage[utf-8]{inputenc}
\usepackage[brazilian]{babel}
\usepackage[top=3cm,left=3cm,bottom=2cm,right=2cm]{geometry}
\usepackage{times}
\usepackage{setspace}
\usepackage{indentfirst}
\usepackage{graphicx}
\usepackage{amsmath}
\usepackage{amssymb}
\usepackage{url}
\usepackage{hyperref}

% Configurar espaçamento duplo
\onehalfspacing

% Configurar a fonte Times New Roman
\renewcommand{\rmdefault}{ptm}

\title{\textbf{Geração de Variáveis Aleatórias}}
\author{Ciência da Computação}
\date{Junho de 2026}

\begin{document}

% CAPA
\begin{titlepage}
\begin{center}
\large
\textbf{PONTIFÍCIA UNIVERSIDADE CATÓLICA DO PARANÁ} \\
Escola Politécnica \\
Curso: Ciência da Computação \\
Disciplina: Métodos Quantitativos \\

\vspace{3cm}

\Large
\textbf{Geração de Variáveis Aleatórias} \\

\vspace{3cm}

\normalsize
Atividade A16 \\
01 de junho de 2026 \\

\vfill

\end{center}
\end{titlepage}

% SUMÁRIO
\tableofcontents
\newpage

% INTRODUÇÃO
\section{Introdução}

Variáveis aleatórias são conceitos fundamentais em probabilidade e estatística. A geração de números aleatórios é essencial em diversas aplicações computacionais, desde simulações científicas até criptografia. Este trabalho apresenta uma pesquisa abrangente sobre técnicas de geração de variáveis aleatórias, métodos para sua avaliação e o funcionamento de geradores modernos.

\newpage

% QUESTÃO 1
\section{Situações que Usam Geradores de Números Aleatórios}

\subsection{Simulação de Sistemas Complexos}

A simulação de Monte Carlo é uma das aplicações mais importantes de geradores de números aleatórios. Utilizando sequências aleatórias, é possível simular comportamentos de sistemas físicos, biológicos e econômicos que seriam difíceis ou impossíveis de resolver analiticamente. Exemplos incluem:

\begin{itemize}
    \item Simulação do comportamento de partículas em física quântica;
    \item Avaliação de risco em portfólios de investimento;
    \item Modelagem de fenômenos naturais como movimentos de fluidos ou propagação de radiação;
    \item Estimativa de integrais definidas através da amostragem aleatória.
\end{itemize}

Na simulação de Monte Carlo, gera-se uma grande quantidade de números aleatórios distribuídos de acordo com probabilidades específicas do sistema. Com suficientes amostras, a distribuição dos resultados converge para valores reais do sistema, permitindo estimativas precisas de grandezas difíceis de calcular diretamente.

\subsection{Criptografia e Segurança}

Geradores de números aleatórios são fundamentais em criptografia moderna. Números verdadeiramente aleatórios (ou pseudo-aleatórios de alta qualidade) são necessários para:

\begin{itemize}
    \item Geração de chaves criptográficas seguras em algoritmos de encriptação como RSA e AES;
    \item Criação de vetores de inicialização (IV) em modos de cifra por bloco;
    \item Geração de nonces e tokens de autenticação únicos;
    \item Salting em funções de hash criptográficas para proteção de senhas.
\end{itemize}

A qualidade do gerador de números aleatórios em aplicações criptográficas é crítica. Um gerador fraco ou previsível comprometeria toda a segurança do sistema. Por isso, aplicações de segurança utilizam geradores criptograficamente seguros, que passam por rigorosos testes estatísticos.

\newpage

% QUESTÃO 2
\section{Método de Geração de Números Aleatórios: Linear Congruente}

\subsection{Descrição do Método}

O Gerador Linear Congruente (LCG) é um dos métodos mais antigos e simples para gerar números pseudo-aleatórios. Sua fórmula é:

\begin{equation}
X_{n+1} = (a \cdot X_n + c) \mod m
\end{equation}

Onde:
\begin{itemize}
    \item $X_n$ é o enésimo número gerado;
    \item $a$ é o multiplicador;
    \item $c$ é o incremento;
    \item $m$ é o módulo;
    \item $X_0$ é a semente inicial.
\end{itemize}

Para gerar números no intervalo $[0, 1)$, normaliza-se o resultado:

\begin{equation}
U_n = \frac{X_n}{m}
\end{equation}

\subsection{Exemplo Prático}

Considere os parâmetros: $a = 1103515245$, $c = 12345$, $m = 2^{31}$ e $X_0 = 1$. A sequência gerada seria:

\begin{align*}
X_1 &= (1103515245 \times 1 + 12345) \bmod 2^{31} = 1103527590 \\
U_1 &= \frac{1103527590}{2^{31}} \approx 0.5140 \\
X_2 &= (1103515245 \times 1103527590 + 12345) \bmod 2^{31} = ... \\
\end{align*}

\subsection{Vantagens e Desvantagens}

\textbf{Vantagens:}
\begin{itemize}
    \item Simples de implementar;
    \item Rápido computacionalmente;
    \item Completamente determinístico e reproduzível.
\end{itemize}

\textbf{Desvantagens:}
\begin{itemize}
    \item Qualidade estatística inferior comparado a métodos modernos;
    \item Períodos curtos com certos parâmetros;
    \item Números sucessivos apresentam correlação (especialmente os bits menos significativos);
    \item Não adequado para aplicações criptográficas;
    \item Padrões de baixa dimensionalidade (problemas com pares e tríadas consecutivas).
\end{itemize}

\newpage

% QUESTÃO 3
\section{Avaliação da Qualidade de Sequências de Números Aleatórios}

\subsection{Testes Estatísticos Básicos}

Para avaliar se uma sequência de números aleatórios é adequada para aplicações, diversos testes estatísticos podem ser aplicados:

\subsubsection{Teste de Uniformidade (Teste de Kolmogorov-Smirnov)}

Verifica se a sequência segue uma distribuição uniforme no intervalo $[0, 1)$. A estatística do teste é:

\begin{equation}
D = \max|F_n(x) - F(x)|
\end{equation}

Onde $F_n(x)$ é a função de distribuição empírica e $F(x)$ é a função de distribuição teórica. Um valor pequeno de $D$ indica boa uniformidade.

\subsubsection{Teste do Qui-Quadrado}

Divide o intervalo $[0, 1)$ em $k$ subintervalos iguais, conta quantos números caem em cada intervalo, e compara com a frequência esperada:

\begin{equation}
\chi^2 = \sum_{i=1}^{k} \frac{(O_i - E_i)^2}{E_i}
\end{equation}

Onde $O_i$ é a frequência observada e $E_i$ é a frequência esperada. A sequência é considerada adequada se $\chi^2$ não exceder o valor crítico da tabela chi-quadrado.

\subsubsection{Teste de Autocorrelação}

Verifica se há dependência entre números consecutivos ou números separados por um lag $\tau$. A autocorrelação é:

\begin{equation}
\rho_\tau = \frac{\sum_{i=1}^{n-\tau}(U_i - \bar{U})(U_{i+\tau} - \bar{U})}{\sum_{i=1}^{n}(U_i - \bar{U})^2}
\end{equation}

Uma sequência de qualidade deve ter $\rho_\tau$ próximo de zero.

\subsection{Testes Estatísticos Avançados}

\subsubsection{Teste de Frequência (Monobit Test)}

Verifica se existem quantidades aproximadamente iguais de zeros e uns na representação binária dos números aleatórios.

\subsubsection{Teste de Runs}

Avalia se há mudanças apropriadas entre valores consecutivos na sequência. Runs muito longos indicam falta de aleatoriedade.

\subsubsection{Teste de Entropia}

Calcula a entropia de Shannon da sequência. Uma sequência verdadeiramente aleatória deve ter entropia máxima para seu conjunto de símbolos.

\subsection{Suites de Teste Conhecidas}

Existem suites de testes padronizadas:

\begin{itemize}
    \item \textbf{Diehard Tests}: 15 testes estatísticos para geradores de números aleatórios;
    \item \textbf{NIST Statistical Test Suite}: Suite desenvolvida pelo NIST com 15 testes para avaliação criptográfica;
    \item \textbf{TestU01}: Biblioteca em C com centenas de testes incluindo os anteriores.
\end{itemize}

\newpage

% QUESTÃO 4
\section{Mersenne Twister}

\subsection{História e Características}

O Mersenne Twister é um gerador de números pseudo-aleatórios desenvolvido por Matsumoto e Nishimura em 1997. É amplamente utilizado em aplicações científicas e em muitas linguagens de programação (Python, Ruby, C++11, etc.). O nome se deve ao fato de usar números primos de Mersenne.

\textbf{Características principais:}
\begin{itemize}
    \item Período muito longo: $2^{19937} - 1$ (aproximadamente $10^{6000}$);
    \item Excelentes propriedades estatísticas de 623 dimensões;
    \item Rápido: gera 32 bits (ou 64 bits na versão MT19937-64) por iteração;
    \item Não é adequado para criptografia, pois é previsível com poucos valores observados;
    \item Implementação simples e bem documentada.
\end{itemize}

\subsection{Algoritmo}

O Mersenne Twister funciona sobre um vetor de estado de 624 valores de 32 bits. O algoritmo pode ser resumido em:

\subsubsection{1. Inicialização (Seeding)}

Dado uma semente $x_0$, inicializa-se o vetor de estado $\mathbf{MT}[0..623]$:

\begin{equation}
\text{MT}[i] = f(i, \text{MT}[i-1])
\end{equation}

Onde $f$ é uma função de inicialização que garante boa distribuição do estado inicial.

\subsubsection{2. Geração de Números}

O processo ocorre em duas fases:

\textbf{Fase 1: Twist (Recorrência Linear)}

A cada iteração, atualiza-se o vetor de estado usando uma operação de recorrência baseada em números primos de Mersenne:

\begin{equation}
\text{MT}[i] = \text{MT}[(i+m) \bmod n] \oplus (\text{upper}(\text{MT}[i]) | \text{lower}(\text{MT}[i+1])) \gg 1) \oplus A
\end{equation}

Onde:
\begin{itemize}
    \item $n = 624$ é o tamanho do vetor de estado;
    \item $m = 397$ é um parâmetro de deslocamento;
    \item $\oplus$ é a operação XOR;
    \item $\gg$ é shift (deslocamento) à direita;
    \item $A$ é uma matriz constante;
    \item \texttt{upper} e \texttt{lower} extraem os bits mais e menos significativos.
\end{itemize}

\textbf{Fase 2: Tempering (Temperagem)}

O valor extraído sofre transformações (tempering) para melhorar propriedades estatísticas:

\begin{align}
y &= \text{MT}[i] \\
y &= y \oplus (y \gg 11) \\
y &= y \oplus ((y \ll 7) \text{ AND } 0x9d2c5680) \\
y &= y \oplus ((y \ll 15) \text{ AND } 0xefc60000) \\
y &= y \oplus (y \gg 18)
\end{align}

O valor $y$ é então retornado como o número pseudo-aleatório.

\subsection{Implementação Conceitual}

O pseudocódigo básico é:

\begin{verbatim}
função inicializa(semente):
    MT[0] = semente
    para i de 1 até 623:
        MT[i] = lowestW bits de (f * (MT[i-1] ^ (MT[i-1] >> (W-2))) + i)

função extrai_número():
    se índice >= 624:
        twist()
    y = MT[índice]
    y = tempering(y)
    índice = índice + 1
    retorna y

função twist():
    para i de 0 até 623:
        x = (MT[i] & UPPER_MASK) | (MT[(i+1) mod 624] & LOWER_MASK)
        MT[i] = MT[(i+397) mod 624] ^ (x >> 1) ^ ((x & 1) * 0x9908b0df)
    índice = 0
\end{verbatim}

\subsection{Vantagens}

\begin{itemize}
    \item Período extremamente longo ($2^{19937} - 1$);
    \item Distribuição uniforme em 623 dimensões;
    \item Rápido e eficiente;
    \item Implementação simples e disponível em muitas linguagens;
    \item Estável e amplamente testado;
    \item Produz sequências de alta qualidade para simulações científicas.
\end{itemize}

\subsection{Limitações}

\begin{itemize}
    \item Não é criptograficamente seguro (pode ser predito com observação de 624 valores consecutivos);
    \item Usa muita memória (624 × 4 bytes = 2.5 KB apenas para estado);
    \item Inicialização simples pode produzir correlações em múltiplas instâncias com sementes próximas;
    \item Alguns bits têm períodos menores que o período total.
\end{itemize}

\subsection{Variantes}

Existem variantes do Mersenne Twister para diferentes propósitos:

\begin{itemize}
    \item \textbf{MT19937-64}: Versão de 64 bits com período $2^{19937} - 1$;
    \item \textbf{SFMT (SIMD-oriented Fast Mersenne Twister)}: Versão otimizada para processadores modernos com instruções SIMD;
    \item \textbf{dSFMT}: Variante que gera números em ponto flutuante diretamente.
\end{itemize}

\newpage

% CONCLUSÃO
\section{Conclusão}

A geração de números aleatórios é uma área fundamental da computação com aplicações vastas em simulação, criptografia e análise estatística. Desde métodos simples como o gerador linear congruente até algoritmos sofisticados como o Mersenne Twister, a evolução desses geradores reflete a importância crescente da qualidade e eficiência em aplicações modernas.

O Mersenne Twister se consolidou como um padrão de fato para aplicações científicas devido ao seu excelente balance entre período longo, qualidade estatística e velocidade de execução. Entretanto, compreender limitações e aplicar testes estatísticos apropriados permanece essencial para garantir que a sequência gerada seja adequada para cada contexto específico.

\newpage

% REFERÊNCIAS
\begin{thebibliography}{99}

\bibitem{wikipedia_va}
Wikipedia. \textbf{Variável Aleatória}. Disponível em: \url{https://pt.wikipedia.org/wiki/Variável_aleatória}. Acesso em: 1 jun. 2026.

\bibitem{wikipedia_mt}
Wikipedia. \textbf{Mersenne Twister}. Disponível em: \url{https://en.wikipedia.org/wiki/Mersenne_Twister}. Acesso em: 1 jun. 2026.

\bibitem{matsumoto1998}
Matsumoto, M.; Nishimura, T. \textbf{Mersenne Twister: a 623-dimensionally equidistributed uniform pseudo-random number generator}. \textit{ACM Transactions on Modeling and Computer Simulation}, v. 8, n. 1, p. 3-30, 1998.

\bibitem{saito2008}
Saito, M.; Matsumoto, M. \textbf{SIMD-oriented Fast Mersenne Twister: a 128-bit Pseudorandom Number Generator}. In: \textit{International Conference on Monte Carlo Methods and Applications}. 2008.

\bibitem{nist2010}
NIST. \textbf{A Statistical Test Suite for Random and Pseudorandom Number Generators for Cryptographic Applications}. Special Publication 800-22 Rev. 3a. Disponível em: \url{https://nvlpubs.nist.gov/nistpubs/Legacy/SP/nistspecialpublication800-22r1a.pdf}. 2010.

\bibitem{knuth1997}
Knuth, D. E. \textbf{The art of computer programming: Volume 2, Seminumerical algorithms}. 3. ed. Addison-Wesley, 1997.

\bibitem{lecuyer2006}
L'Ecuyer, P. \textbf{Handbook of Computational Statistics: Concepts and Methods}. In: \textit{Pseudo-Random Number Generators}. Springer, 2006. Cap. 3.

\bibitem{nvidia2010}
NVIDIA. \textbf{Mersenne Twister Applied to Surface Scattering}. CUDA SDK. 2010.

\end{thebibliography}

\end{document}
